\section{Spectral sufficiency}
\label{sec:sufficiency}

Fix $n$ and consider the observed matrix $X_n\in\Sym_n$. Recall from
Section~\ref{sec:notation} that $D(X)\in\R^n$ is the vector of eigenvalues
in decreasing order and, on the event of simple spectrum (which has
probability one under every rotatable law with $\sigma>0$), $V(X)\in\OO(n)$
is the sign-normalized eigenvector matrix, so
$X = V\,\diag(D)\,V^\top$.

\begin{theorem}[Spectral sufficiency]
\label{thm:suff}
Let $\mathcal{P}$ be any family of orthogonally invariant laws on
$\Sym_n$, for example the fully observed $n\times n$ corners, diagonal
retained, of the jointly rotatable laws $R_{\rho,\sigma,\blam}$ of
Theorem~\ref{thm:rep} and arbitrary mixtures of them; the hollow laws
$P^{\mathrm{net}}_{\sigma,\blam}$ are not orthogonally invariant
(Remark~\ref{rem:hollow}) and are excluded. Then:
\begin{enumerate}[label=(\roman*),itemsep=1pt,topsep=2pt]
\item For every $P\in\mathcal{P}$, conditionally on $D(X)$ the matrix $X$
is distributed as $U\,\diag(D)\,U^\top$ with $U\sim\Haar_n$
independent of $D$; in particular the conditional law of $X$ given $D$ is
the same for all $P\in\mathcal{P}$.
\item $D(X)$ is sufficient for $\mathcal{P}$ in the strong,
common-kernel sense: one Markov kernel, free of $P$, is a regular
conditional law of $X$ given $D(X)$ under every $P\in\mathcal P$; no
domination assumption on $\mathcal{P}$ is required.
\item On the simple-spectrum event, the frame $V(X)$ follows the
sign-normalized Haar law $\nu_n$, the pushforward of
$\Haar_n$ under columnwise sign normalization, which is not
Haar measure on $\OO(n)$, independently of $D(X)$ and identically for
every $P\in\mathcal{P}$: the frame is an ancillary complement to the
spectrum, and $(D(X),V(X))$ recovers $X$.
\item Conversely, if a law $P$ on $\Sym_n$ satisfies (i) then it is
orthogonally invariant: spectral sufficiency with Haar conditionals
characterizes rotatable models.
\end{enumerate}
\end{theorem}

The proof is a disintegration argument, given in the Supplement with
an explicit measurable construction of the eigenframe $V$; the value
lies in practice, combined with Theorem~\ref{thm:rep}: the eigenvalues
carry all information about every parameter of every rotatable model,
in the exact finite-sample sense of sufficiency, and
$X\leftrightarrow(D,V)$ splits the data into a sufficient part and an ancillary complement with exactly known
law.

The consequence for network statistics is Corollary~S1 in the
Supplement. For any measurable network statistic $S:\Sym_n\to\R^q$
(clustering, modularity, path length, small-worldness, strengths,
motifs) and any family of rotatable laws: the conditional law of
$S(X)$ given $D(X)$ is parameter-free and exactly simulable by Haar
rotation of the observed spectrum; $\E[S(X)\mid D]$ is a spectral
statistic and any procedure based on $S$ is weakly dominated by one
based on $D$ (Rao--Blackwell); and a scale-invariant statistic of the
frame $V(X)$ alone is ancillary. Without the invariance
\eqref{eq:rotdef} every part of this fails: part (iv) of
Theorem~\ref{thm:suff} shows the Haar conditional characterizes
rotatability, so for any non-rotatable law the frame carries
information and the reduction to $D$ discards it.

In the language of the motivating applications: under a rotatable
law, the conditional distribution of modularity or clustering given
the eigenvalues is that of a uniformly random rotation of those
eigenvalues, whatever the underlying $(\sigma,\blam)$; this is
falsifiable, Section~\ref{sec:test} builds the exact test, and
Section~\ref{sec:experiments} carries it out. Two remarks fix the
scope of the claim.

\begin{remark}[What the claim does and does not say]
\label{rem:scope}
Corollary~S1 does not say that non-spectral summaries are
marginally uninformative: under an ergodic law, clustering has a
nondegenerate limit determined by $(\sigma,\blam)$ and will correlate with
any phenotype the spectrum predicts; in the cohort experiment of
Section~\ref{sec:experiments}, classifiers using non-spectral summaries
alone perform well when the truth is rotatable. The claim is conditional:
given $D$, the summaries add exactly nothing under the null. Whether this
conditional statement approximately describes real networks is an empirical
question examined in Section~\ref{sec:experiments}.
\end{remark}

\begin{remark}[Extremal-family perspective]
Theorem~\ref{thm:suff} is an instance of the symmetry--sufficiency
duality for extremal families
\citep{lauritzen1988extremal,diaconis1984partial}: the ergodic
decomposition presents the rotatable networks as an extremal family
whose canonical sufficient statistic is the spectrum, in analogy with
the empirical measure for exchangeable sequences; under exchangeability
alone no comparable reduction exists.
\end{remark}
